% ch07_u03 —— 原书 p.256--p.265（7.4.11 -- 7.6.10）
% 本单元起于 p.256 首段（承接 p.255 的 7.4.10），末尾停在 7.6.10 的中途；尾部接 p.266。

\SourcePage{256}

我们用对偶来定义 DG 函子 $t_{-}:C\sec_{-}\to C\mathrm{hot}_{-}$ 与 $s_{+}:C\mathrm{hot}_{-}\to C\sec_{+}$。注意，复合 $s_{+}t_{-}:C\sec_{-}\to C\sec_{+}$ 与 7.4.2 中的函子 $c_{+}$ 一致；类似地，$s_{-}t_{+}=c_{-}$。函子 $t_{-}s_{-}:C\mathrm{hot}_{+}\to C\mathrm{hot}_{-}$ 与 $t_{+}s_{+}:C\mathrm{hot}_{-}\to C\mathrm{hot}_{+}$ 互为伴随（正如 7.4.2 中的函子 $c_{\pm}$）。

设 $f:A\to B$ 是 $C\mathrm{hot}_{\pm}$ 的同伦范畴 $K\mathrm{hot}_{\pm}$ 中的态射。若所有态射 $f_{m}:=f_{\mathrm{id}_{\Delta_{m}}}:A_{m}\to B_{m}$ 都是拟同构，则称 $f$ 为\emph{拟同构}。拟同构构成局部化族。相应的局部化三角范畴记为 $D\mathrm{hot}_{\pm}$。

函子 $s_{\pm}$、$t_{\pm}$ 保持拟同构，因此给出导出范畴之间的函子。这些函子复合的伴随态射都是拟同构。于是导出范畴 $D\sec_{\pm}$ 与 $D\mathrm{hot}_{\pm}$ 典范地等同。

% [Tutor Note: 原书 p.256 此处拼作「canonicaly」（应为 canonically），底本如此。]

\emph{注.} (i) 复形 $A\in D\mathrm{hot}_{+}$ 属于 $D_{\mathrm{tot}}$，当且仅当对任意 $\alpha:\Delta_{m}\to\Delta_{m+1}$，其 $\alpha$-分量 $d_{A\alpha}:\alpha_{\cdot}A_{m}\to A_{m+1}$ 是复形的拟同构（$A_{m}$ 上的微分为 $d_{A\,\mathrm{id}_{\Delta_{m}}}$，$A_{m+1}$ 亦然）。

(ii) 若范畴 $\mathcal{M}_{n}$ 有许多内射对象，则 $K^{+}\mathrm{hot}_{+}$ 有许多内射对象。对偶地，若 $\mathcal{M}_{n}$ 有许多投射对象，则 $K^{-}\mathrm{hot}_{-}$ 有许多投射对象（参看 7.4.9）。

\bdp{7.4.11} 上述某些构造在下面稍一般的设定中仍有意义。考虑任一在 $(\Delta)$ 上余纤维化的 DG 范畴族 $\mathcal{C}_{\cdot}$。有 DG 范畴 $\mathcal{C}\sec_{\pm}=\sec_{\pm}(\mathcal{C}_{\cdot})$（其定义与 7.4.2 中的范畴 $\sec_{\pm}(\mathcal{M}_{\cdot})$ 完全相同），以及相应的同伦范畴。按 7.4.2 的方式，在 $\pm$ 两类范畴之间定义伴随函子 $c_{\pm}$。

再设我们有 7.4.1 中的 $\mathcal{M}$，以及一族与纤维范畴结构相容的上同调函子 $H:\mathcal{C}_{\cdot}\to\mathcal{M}_{\cdot}$。于是得到相应的上同调函子 $H_{\pm}:\mathcal{C}\sec_{\pm}\to\sec_{\pm}$。把同伦范畴对 $H$-拟同构局部化，得到导出范畴 $D\sec_{\pm}$。如引理 7.4.4 那样，函子 $c_{\pm}$ 把范畴 $D\sec_{\pm}$ 等同起来，故可简记为 $D\sec$。接着定义\SourcePage{257}范畴 $\mathcal{C}_{\mathrm{tot}\pm}$ 等，如 7.4.5 所述。引理 7.4.6 仍然成立，故有全三角子范畴 $D_{\mathrm{tot}}\subset D\sec$，以及上同调函子 $H:D_{\mathrm{tot}}\to\mathcal{M}_{\mathrm{tot}}$。

\bdsec{7.5. 光滑叠上的 $\mathcal{D}$-模理论 II}

\bdp{7.5.1} 设 $\mathcal{Y}$ 是任意的光滑代数叠。设 $U_{\cdot}$ 是 $\mathcal{Y}$ 的超覆叠，使得每个 $U_{n}$ 都是（光滑）拟紧分离代数空间的无交并（例如仿射概形）。称这样的 $U_{\cdot}$ 为\emph{容许}超覆叠。把 $U_{\cdot}$ 看作 $(\Delta)^{\circ}$-代数空间。范畴 $\mathcal{M}(U_{\cdot})$ 构成 7.4.1 中的 $(\Delta)$-族阿贝尔范畴；相应的范畴 $\mathcal{M}_{\mathrm{tot}}$ 就是 $\mathcal{M}(\mathcal{Y})$。依 7.4.7，得到相应的 t-范畴 $D_{\mathrm{tot}}=D_{\mathrm{tot}}(U_{\cdot},\mathcal{D})$，其核心为 $\mathcal{M}(\mathcal{Y})$。

我们也可以考虑 DG 范畴 $C(U_{\cdot},\Omega)$，连同上同调函子 $H_{\mathcal{D}\cdot}:C(U_{\cdot},\Omega)\to\mathcal{M}(U_{\cdot})$，其中 $H_{\mathcal{D}n}F_{n}=H_{\mathcal{D}}F_{n}[\dim U_{n}/\mathcal{Y}]$，$F_{n}\in C(U_{n},\Omega)$；然后应用 7.4.11。得到三角范畴 $D_{\mathrm{tot}}(U_{\cdot},\Omega)$，以及上同调函子 $H_{\mathcal{D}}:D_{\mathrm{tot}}(U_{\cdot},\Omega)\to\mathcal{M}(\mathcal{Y})$。

范畴 $D_{\mathrm{tot}}(U_{\cdot},\mathcal{D})$ 与 $D_{\mathrm{tot}}(U_{\cdot},\Omega)$ 典范地等同。确切地说，有函子 $\Omega_{\cdot}:C(U_{\cdot},\mathcal{D})\to C(U_{\cdot},\Omega)$，$\Omega_{n}(M_{n}):=\Omega M_{n}[-\dim U_{n}/\mathcal{Y}]$。该函子与 DG 结构、纤维范畴结构以及与上同调函子相容（即 $H=H_{\mathcal{D}\cdot}\Omega_{\cdot}$）。因此它给出正合函子
\begin{equation}\tag{315}
\Omega:D_{\mathrm{tot}}(U_{\cdot},\mathcal{D})\to D_{\mathrm{tot}}(U_{\cdot},\Omega).
\end{equation}

该函子是范畴等价。事实上，尽管 $C(U_{\cdot},\Omega)$ 与 $C(U_{\cdot},\mathcal{D})$ 之间的函子 $\mathcal{D}$ \emph{不}与纤维范畴结构相容，它仍给出函子 $\mathcal{D}:C\sec_{-}(U_{\cdot},\Omega)\to C\sec_{-}(U_{\cdot},\mathcal{D})$，$(\mathcal{D}F)_{n}=\mathcal{D}F_{n}[\dim U_{n}/\mathcal{Y}]$（用 7.2.8 来定义各个 $\alpha^{*}$）。这个 $\mathcal{D}$ 是相应 $\Omega$ 函子的左伴随，且与上同调函子相容。$\mathcal{D}$-$\Omega$ 伴随态射都是拟同构（见 7.2.4、7.2.5），故 $\mathcal{D}$ 给出 (315) 的逆函子。

\SourcePage{258}

我们把这样等同起来的范畴 $D_{\mathrm{tot}}(U_{\cdot},\mathcal{D})$ 与 $D_{\mathrm{tot}}(U_{\cdot},\Omega)$ 简记为 $D_{\mathrm{tot}}(U_{\cdot})$。

\bdp{7.5.2} \emph{命题.} 对 $\mathcal{Y}$ 的不同容许覆盖，t-范畴 $D_{\mathrm{tot}}(U_{\cdot})$ 之间存在典范等同。

证明见下文 7.5.5。把这样等同起来的范畴记为 $D(\mathcal{Y})$；它是核心为 $\mathcal{M}(\mathcal{Y})$ 的 t-范畴。

在证明 7.5.2 之前，先说明：若 $\mathcal{Y}$ 满足条件 (310)，则确实得到与 7.3.2 相同的范畴 $D(\mathcal{Y})$。顺便，这蕴含 7.3.4。

选取 $\mathcal{Y}$ 的超覆叠 $U_{\cdot}$，使各 $U_{n}$ 为仿射概形。有显然的正合函子（限制到 $U_{\cdot}$）
\begin{equation}\tag{316}
r:D(\mathcal{Y},\Omega)\to D_{\mathrm{tot}}(U_{\cdot},\Omega).
\end{equation}

\bdp{7.5.3} \emph{引理.} 函子 $r$ 是范畴等价。

\emph{证明.} 我们来构造逆函子。对 $F\in K_{\mathrm{tot}+}(\Omega)$，用各项为 $\pi_{\cdot}F^{ab}:=\pi_{b}(F^{a})$ 的 Čech 双复形的全复形，定义 $\mathcal{Y}$ 上的 $\Omega$-复形 $\pi_{\cdot}F$；于是 $(\pi_{\cdot}F)^{n}=\bigoplus_{a+b=n}F^{ab}$，这里 $\pi_{b}$ 是投影 $U_{b}\to\mathcal{Y}$。这样得到正合函子 $\pi_{\cdot}:K_{\mathrm{tot}+}(\Omega)\to K(\mathcal{Y},\Omega)$。该函子保持 $\mathcal{D}$-拟同构（因为由 (310)，投影 $\pi_{b}$ 是仿射的），故给出函子 $D_{\mathrm{tot}}(U_{\cdot},\Omega)\to D(\mathcal{Y},\Omega)$。

验证这个函子是 $r$ 的逆，留给读者（提示：对如上 $F$，伴随拟同构 $\pi'_{\Omega}\pi_{\cdot}F\to F$ 来自 $C\sec_{-}(U_{\cdot},\Omega)$ 中的典范态射 $\pi'_{\Omega}\pi_{\cdot}F\to c_{-}F$）。\hfill$\Box$

\bdp{7.5.4} \emph{注.} 上述引理可按如下方式移植到 7.3.12 的设定中。设 $\mathcal{Y}$ 是任意的光滑叠，其对角态射 $\mathcal{Y}\to\mathcal{Y}\times\mathcal{Y}$ 为仿射。则 7.3.12 与 7.5.1 中定义的范畴 $D(\mathcal{Y})$ 典范等价。事实上，设 $\mathcal{Y}_{i}$ 是 7.3.12 中 $\mathcal{Y}$ 的一列开子叠，$V_{i}\to\mathcal{Y}$ 是覆盖，各 $V_{i}$ 为仿射概形。\SourcePage{259}于是各 $V_{i}$ 构成 $\mathcal{Y}$ 的覆盖。设 $U_{\cdot}$ 是相应的 Čech 超覆叠。因此 $U_{a}$ 是分支 $U_{\alpha}$ 的无交并，这些分支由序列 $\alpha=(\alpha_{1},\alpha_{2},\ldots)$ 标号，$\alpha_{i}\ge0$，$\Sigma\alpha_{i}=a+1$；其中 $U_{\alpha}$ 是 $\alpha_{1}$ 个 $V_{1}$、$\alpha_{2}$ 个 $V_{2}$……在 $\mathcal{Y}$ 上的纤维积。对 $F\in C(\mathcal{Y}_{\cdot},\Omega)$，置 $F_{U\alpha}:=F_{i_{\alpha}U_{\alpha}}$，其中 $i_{\alpha}$ 是使 $\alpha_{i}$ 非零的最小 $i$（注意 $U_{\alpha}\in\mathcal{Y}_{sm}$。这些 $F_{U_{\alpha}}$ 按显然的方式构成 $U_{\cdot}$ 上的 $\Omega$-复形 $F$，它属于 $C_{\mathrm{tot}}(U_{\cdot},\Omega)$。

% [Tutor Note: 原书 p.259「(note that U_α ∈ Y_sm.」左括号未闭合，底本如此，译文依原样。]

函子 $C(\mathcal{Y}_{\cdot},\Omega)\to C_{\mathrm{tot}}(U_{\cdot},\Omega)$ 与函子 $H_{\mathcal{D}}$ 交换，故给出三角函子
\begin{equation}\tag{317}
r:D(\mathcal{Y}_{\cdot},\Omega)\to D_{\mathrm{tot}}(U_{\cdot},\Omega).
\end{equation}

验证该函子是范畴等价，以及 7.3.12 与 7.5.2 意义下 $D(\mathcal{Y})$ 的相应等同不依赖 $\mathcal{Y}_{\cdot}$ 与 $V_{\cdot}$ 这些辅助数据，均留给读者。

\bdp{7.5.5} \emph{7.5.2 的证明.} 我们需要对不同的 $U_{\cdot}$ 典范地等同 t-范畴 $D_{\mathrm{tot}}(U_{\cdot})$。设 $U'_{\cdot}$ 是另一个容许超覆叠。首先用某些辅助数据 $V$ 来定义 t-正合函子 $\Phi=\Phi_{V}:D_{\mathrm{tot}}(U_{\cdot})\to D_{\mathrm{tot}}(U'_{\cdot})$。然后证明 $\Phi$ 实际上不依赖 $V$，且是范畴等价。

我们的 $V$ 是 $\mathcal{Y}$ 上的 $(\Delta)^{\circ}\times(\Delta)^{\circ}$-代数空间 $V_{\cdot\cdot}$，连同光滑态射 $\pi:V_{mn}\to U_{m}$、$\pi':V_{mn}\to U'_{n}$。假定 $\pi$、$\pi'$ 按显然的方式与 $(\Delta)$ 投影相容，$\pi'_{\cdot n}:V_{\cdot n}\to U'_{n}$ 是超覆叠，$\pi'_{mn}:V_{mn}\to U'_{n}$ 是仿射态射。

对 $F\in K_{\mathrm{tot}+}(U_{\cdot},\Omega)$，有 $\Omega$-复形 $F_{V_{n}}\in K_{\mathrm{tot}+}(V_{\cdot n},\Omega)$，即 $F$ 到 $V_{\cdot n}$ 的拉回。置 $\Phi_{V_{n}}F:=\pi'_{\cdot}F_{V_{n}}$（记号见 7.5.3 的证明）。这是 $U'_{n}$ 上的 $\Omega$-复形。这些 $\Omega$-复形 $\Phi_{V_{n}}F$ 按显然的方式构成 $\Omega$-复形 $\Phi_{V}F\in K_{\mathrm{tot}+}(U'_{\cdot},\Omega)$，并满足 $H_{\mathcal{D}}F=H_{\mathcal{D}}\Phi_{V}F$。因此有 t-正合函子 $\Phi_{V}:D_{\mathrm{tot}}(U_{\cdot},\Omega)\to D_{\mathrm{tot}}(U'_{\cdot},\Omega)$，它在核心 $\mathcal{M}(\mathcal{Y})$ 之间诱导恒等函子。

设有如上 $V_{1}$ 与 $V_{2}$。为等同函子 $\Phi_{V_{i}}$，选取另一个如上 $V$，连同嵌入 $V_{1},V_{2}\subset V$\SourcePage{260}，且这些嵌入与所有投影相容，这些投影把 $(V_{1})_{mn}$、$(V_{2})_{mn}$ 与 $V_{mn}$ 的若干连通分支之并等同。这些嵌入诱导投影 $\Phi_{V}F\to\Phi_{V_{1}}$、$\Phi_{V}F\to\Phi_{V_{2}}F$，它们显然是拟同构。

% [Tutor Note: 原书 p.260 第一个投影右端作「Φ_V F → Φ_{V_1}」，末端缺 F，底本如此，译文依原样。]

于是我们等同了导出范畴之间的函子 $\Phi_{V_{i}}$。验证该等同不依赖辅助数据 $V$，留给读者。

于是有典范函子 $\Phi=\Phi_{UU'}:D_{\mathrm{tot}}(U_{\cdot},\Omega)\to D_{\mathrm{tot}}(U'_{\cdot},\Omega)$。若 $U''_{\cdot}$ 是第三个超覆叠，则有函子的典范同构 $\Phi_{UU''}=\Phi_{U'U''}\Phi_{UU'}$；其定义以及通常相容性的验证，都留给读者。由于 $\Phi_{UU}$ 是恒等函子，可见这些 $\Phi$ 同时等同了所有范畴 $D_{\mathrm{tot}}(U_{\cdot})$。\hfill$\Box$

\bdp{7.5.6} 设 $f:\mathcal{Y}\to\mathcal{Z}$ 是光滑叠的拟紧态射。我们来定义前推函子 $f_{*}:D(\mathcal{Y})^{+}\to D(\mathcal{Z})^{+}$。为此，考虑 $\mathcal{Y}$ 的任一容许超覆叠 $U_{\cdot}$ 与 $\mathcal{Z}$ 的任一容许超覆叠 $W_{\cdot}$。得到 $(\Delta)^{\circ}\times(\Delta)^{\circ}$-代数空间 $U_{\cdot}\times_{\mathcal{Z}}W_{\cdot}$。可以找到 $(\Delta)^{\circ}\times(\Delta)^{\circ}$-代数空间 $V_{\cdot\cdot}$，连同态射 $\phi=(\phi_{1},\phi_{2}):V_{\cdot\cdot}\to U_{\cdot}\times_{\mathcal{Z}}W_{\cdot}$，使得投影 $V_{mn}\to U_{m}$ 光滑，$V_{mn}\to W_{n}$ 仿射，且 $V_{\cdot n}\to\mathcal{Y}\times_{\mathcal{Z}}W_{n}$ 是超覆叠。现在对 $F\in K^{+}_{\mathrm{tot}+}(U_{\cdot},\Omega)$，令 $F_{V_{n}}\in K^{+}_{\mathrm{tot}+}(V_{\cdot n},\Omega)$ 是它到 $V_{\cdot n}$ 的拉回。用各项为 $\phi_{2\cdot}F_{V_{n}}$ 的 Čech 双复形的全复形，定义 $W_{n}$ 上的 $\Omega$-复形 $f_{\cdot}F_{n}$。

% [Tutor Note: 原书 p.260 上一行作 F^{V_n}，此处作 F_{V_n}（下标），指同一对象；底本如此，译文依原样。]

这些 $\Omega$-复形构成 $K^{+}_{\mathrm{tot}+}(W_{\cdot},\Omega)$ 的对象 $f_{\cdot}F$。函子 $f_{\cdot}:K^{+}_{\mathrm{tot}+}(U_{\cdot},\Omega)\to K^{+}_{\mathrm{tot}}(W_{\cdot},\Omega)$ 保持 $\mathcal{D}$-拟同构，因此给出函子 $f_{*}:D(\mathcal{Y})^{+}\to D(\mathcal{Z})^{+}$。验证 $f_{*}$ 的构造不依赖 $U$、$W$、$V$ 的辅助选取，且与 $f$ 的复合相容，留给读者。

% [Tutor Note: 原书 p.260 第二个范畴作 K^+_{tot}(W_·,Ω)，漏了下标 +，底本如此，译文依原样。]

光滑叠的光滑态射 $f:\mathcal{Y}\to\mathcal{Z}$ 给出 t-正合函子 $f^{\dagger}=f^{\dagger}_{\Omega}:D(\mathcal{Z})\to D(\mathcal{Y})$。具体地说，选取 $\mathcal{Y}$ 的容许超覆叠 $U_{\cdot}$、$\mathcal{Z}$ 的容许超覆叠 $W_{\cdot}$，以及与 $f$ 相容的态射 $f_{\cdot}:U_{\cdot}\to W_{\cdot}$。函子 $f^{\cdot}_{\Omega}:K_{\mathrm{tot}\pm}(W_{\cdot},\Omega)\to K_{\mathrm{tot}\pm}(U_{\cdot},\Omega)$ 保持 $\mathcal{D}$-拟同构，故给出导出范畴之间的函子 $f^{\cdot}_{\Omega}$。我们把这个定义的检验留给读者：\SourcePage{261}它不依赖辅助选取，我们的拉回函子与 $f$ 的复合相容，且当 $f$ 拟紧时函子 $f^{\cdot}_{\Omega}$ 是 $f_{*}$ 的左伴随。

\bdp{7.5.7}

% [Tutor Note: 原书 p.261 此编号段后无正文（仅印出编号），底本如此。]

\bdp{7.5.8} \emph{注.} (i) 也可以尝试用适当的非拟凝聚 $\Omega$-复形来定义 $D(\mathcal{Y})$，方式类似于 [LMB93]6.3 中 $\mathcal{O}$-模导出范畴的定义。这样的定义很可能给出同一个范畴 $D^{+}(\mathcal{Y})$。

(ii) 导出范畴的「局部」构造在 $\mathcal{O}$-模的设定中也很方便。例如，它有助于把代数叠的余切复形定义为导出范畴中的真正对象（而不只是像 [LMB93]9.2 那样取其截断的投射极限），也有助于处理 Grothendieck-Serre 对偶。

(iii) 把 $\mathcal{D}$-模换成反常层，就得到局部有限型代数叠上可构造层导出范畴的一个方便定义。

\bdsec{7.6. 等变设定}

\bdp{7.6.1} 我们来解释（有限维）Hecke 模式的 7.1.1 (a)、(b) 两部分。设 $G$ 是代数群，$K\subset G$ 是代数子群。为简单起见，假定 $K$ 仿射；此时下面的叠满足 7.3.1 的条件 (310)。置$^{*)}$\ $\mathcal{H}^{c}:=C(K\backslash G/K,\Omega)$，$\mathcal{H}:=D(K\backslash G/K)$。\footnote{$^{*)}$此处上标「c」表示我们处理的是真正的 DG 复形范畴，而非导出范畴。}分别称这两个范畴为\emph{预 Hecke} 范畴与 \emph{Hecke} 范畴。它们带有如下定义的典范幺半结构。

考虑叠的态射
\begin{equation}\tag{318}
(K\backslash G/K)\times(K\backslash G/K)\xleftarrow{p}K\backslash G\times_{K}G/K\xrightarrow{\bar m}K\backslash G/K.
\end{equation}

\SourcePage{262}

这里 $G\times_{K}G$ 是 $G\times G$ 模去 $K$-作用 $k(g_{1},g_{2})=(g_{1}k^{-1},kg_{2})$ 的商，$p$ 是显然的投影，$\bar m$ 是乘法映射。对 $F_{1},F_{2}\in\mathcal{H}^{c}$，置 $F_{1}\overset{c}{\circledast}F_{2}:=\bar m_{\cdot}p^{\cdot}_{\Omega}(F_{1}\boxtimes F_{2})$。\emph{卷积张量积} $\overset{c}{\circledast}$ 满足显然的结合约束，故 $\mathcal{H}^{c}$ 上有幺半结构。把卷积张量积 $\circledast:\mathcal{H}\times\mathcal{H}\longrightarrow\mathcal{H}$ 定义为 $\overset{c}{\circledast}$ 的右导出函子。有 $F_{1}\circledast F_{2}=\bar m_{*\cdot}p^{\cdot}_{\Omega}(F_{1}\boxtimes F_{2})$；若 $\Omega$-复形 $F_{1}$、$F_{2}$ 是松的（见 7.3.7），则 $F_{1}\circledast F_{2}=F_{1}\overset{c}{\circledast}F_{2}$。于是 $\circledast$ 的结合约束可由 $\overset{c}{\circledast}$ 的结合约束推出，故 $\mathcal{H}$ 是幺半三角范畴。$\mathcal{H}^{c}$ 与 $\mathcal{H}$ 有单位对象 $E$：有 $E_{G}=i_{K\cdot}\Omega_{K}$（这里 $i_{K}:K\hookrightarrow G$ 是嵌入）。

设 $Y$ 是带 $G$-作用的光滑簇。考虑叠 $\mathcal{B}:=K\backslash Y$。$\mathcal{D}(\mathcal{B})$ 上的 Hecke 作用来自图表
\begin{equation}\tag{319}
(K\backslash G/K)\times\mathcal{B}\xleftarrow{p_{Y}}K\backslash(G\times_{K}Y)\xrightarrow{\bar m_{Y}}\mathcal{B}.
\end{equation}

即：对 $F\in\mathcal{H}^{c}$、$T\in C(\mathcal{B},\Omega)$，置 $F\overset{c}{\circledast}T:=\bar m_{Y\cdot}p^{\cdot}_{Y\Omega}(F\boxtimes T)$。与上面一样，$\overset{c}{\circledast}$ 满足显然的结合约束，故 $C(\mathcal{B},\Omega)$ 是幺半 $\mathcal{H}^{c}$-模。把 $\circledast:\mathcal{H}\times D(\mathcal{B})\longrightarrow D(\mathcal{B})$ 定义为 $\overset{c}{\circledast}$ 的右导出函子。有 $F\circledast T=\bar m_{Y*\cdot}p^{\cdot}_{Y\Omega}(F\boxtimes T)$；若 $F$、$T$ 松（见 7.3.7），则 $F\circledast T=F\overset{c}{\circledast}T$。因此 $D(\mathcal{B})$ 是 $\mathcal{H}$-模。

\bdp{7.6.2} \emph{注.} (i) 上面的定义中之所以能考虑无界导出范畴，是因为投影 $\bar m$、$\bar m_{Y}$ 可表（见 7.3.11(ii)）。

(ii) 若 $f:Z\to Y$ 是带 $G$-作用的光滑簇的态射，则 $f_{*}:D(K\backslash Z)\to D(K\backslash Y)$ 是 $\mathcal{H}$-模的态射。

\bdp{7.6.3} 设 $Y$ 是带仿射代数群 $K$ 作用的光滑簇。考虑叠 $\mathcal{B}:=K\backslash Y$。在 7.6 的其余部分，我们将用 $Y$ 上适当的等变复形来描述 $D(\mathcal{B})$。我们还将引入某个导出范畴 $D(K^{\backslash}_{\backslash}Y)$，它介于 $D(K\backslash Y)$ 与 $D(Y)$ 之间，在 7.7 中会用到。

\SourcePage{263}

置 $K_{\Omega}=(K,\Omega_{K})$，$K^{\cdot}_{\Omega}=(K,\Omega^{\cdot}_{K})$（于是 $K^{\cdot}_{\Omega}$ 就是略去 de Rham 微分后的 $K_{\Omega}$）。它们分别是 DG 环化空间范畴与分次环化空间范畴中的群对象。记 $\mathfrak{k}$、$\mathfrak{k}_{\Omega}$、$\mathfrak{k}^{\cdot}_{\Omega}$ 分别为 $K$、$K_{\Omega}$、$K^{\cdot}_{\Omega}$ 的李代数。作为普通复形，$\mathfrak{k}_{\Omega}$ 等于 $\mathrm{id}_{\mathfrak{k}}$ 的锥，故 $\mathfrak{k}^{0}_{\Omega}=\mathfrak{k}=\mathfrak{k}^{-1}_{\Omega}$。由于 $K$ 是 $K_{\Omega}$ 与 $K^{\cdot}_{\Omega}$ 的子群，我们得到相应的 Harish-Chandra 对 $(\mathfrak{k}_{\Omega},K)$、$(\mathfrak{k}^{\cdot}_{\Omega},K)$。注意，$K_{\Omega}$-模与 DG $(\mathfrak{k}_{\Omega},K)$-模相同，$K^{\cdot}_{\Omega}$-模与分次 $(\mathfrak{k}^{\cdot}_{\Omega},K)$-模相同。

% [Tutor Note: 原书 p.263「the category of DG ringed spaced」多一字母 d（应为 spaces），底本如此。]

$Y$ 上的 $K$-作用给出 $K_{\Omega}$ 在 $Y_{\Omega}=(Y,\Omega)$ 上的作用，从而给出 $K^{\cdot}_{\Omega}$ 在 $Y^{\cdot}_{\Omega}=(Y,\Omega^{\cdot})$ 上的作用。对分次 $\Omega^{\cdot}_{Y}$-模 $F^{\cdot}_{Y}$，$F^{\cdot}_{Y}$ 上的 $K^{\cdot}_{\Omega}$-作用等同于 $(\mathfrak{k}^{\cdot}_{\Omega},K)$-作用。具体地说，它是 $F^{\cdot}_{Y}$ 上的 $K$-作用，连同 $K$-等变态射 $\mathfrak{k}\otimes F^{\cdot}_{Y}\to F^{\cdot-1}_{Y}$，$\xi\otimes f\mapsto i_{\xi}(f)$（假定 $K$ 以伴随方式作用在 $\mathfrak{k}$ 上），使得对任意 $\xi\in\mathfrak{k}$ 与 $\nu\in\Omega^{1}_{Y}$ 有 $i_{\xi}(\nu f)=\langle\xi,\nu\rangle f+\nu i_{\xi}(f)$，$i^{2}_{\xi}=0$。

\bdp{7.6.4} 设 $F_{Y}$ 是 $Y$ 上的 $\Omega$-复形。$F_{Y}$ 上的 \emph{$K$-作用} 是分次 $\mathcal{O}_{Y}$-模 $F^{\cdot}_{Y}$ 上的 $K$-作用，使得对任意 $k\in K$，平移 $k^{*}F^{\cdot}_{Y}\xrightarrow{\sim}F^{\cdot}_{Y}$ 是 $\Omega$-复形的态射（即它与微分交换）。$F_{Y}$ 上的 \emph{$K_{\Omega}$-作用} 是把 $F_{Y}$ 看作 $Y_{\Omega}$ 上的 DG 模时 $K_{\Omega}$ 的作用。换言之，它是分次 $\Omega^{\cdot}_{Y}$-模 $F^{\cdot}_{Y}$ 上的 $K^{\cdot}_{\Omega}$-作用，使得 $K$ 按 $\Omega$-复形作用于 $F_{Y}$，而 $\mathfrak{k}_{\Omega}$ 按 DG 李代数作用于 $F_{Y}$。后一条件意味着对任意 $\xi\in\mathfrak{k}$ 有 $di_{\xi}+i_{\xi}d=\mathrm{Lie}_{\xi}$（这里 $\mathrm{Lie}$ 是由 $K$-作用得到的 $\mathfrak{k}$ 在 $F^{\cdot}_{Y}$ 上的作用）。带 $K$-作用的 $\Omega$-复形称为\emph{弱 $K$-等变} $\Omega$-复形，带 $K_{\Omega}$-作用的称为 \emph{$K_{\Omega}$-等变} $\Omega$-复形。

显然，对叠 $\mathcal{B}:=K\backslash Y$ 上的任意 $\Omega$-复形 $F$，$\Omega$-复形 $F_{Y}$ 自动带有 $K_{\Omega}$-作用。

\bdp{7.6.5} \emph{引理.} 函子 $C(K^{\backslash}_{\backslash}Y,\Omega)\longrightarrow$（$Y$ 上的 $K_{\Omega}$-等变 $\Omega$-复形）是 DG 范畴的等价。\hfill$\Box$

\SourcePage{264}

\bdp{7.6.6} \emph{注.} 假定处于 7.6.1 的情形。设 $m:K\backslash G\times G/K\to K\backslash G/K$ 是乘法映射。置 $F_{1}\widetilde{\circledast}F_{2}=m_{\cdot}(F_{1K\backslash G}\boxtimes F_{2G/K})$；这是 $K\backslash G/K$ 上的 $\Omega$-复形。沿投影 $G\times G\to G\times_{K}G$ 的纤维的 $K$-作用给出 $F_{1}\widetilde{\circledast}F_{2}$ 上的 $K_{\Omega}$-作用（取 $K\backslash G/K$ 上的平凡 $K$-作用）。其不变量与 $F_{1}\overset{c}{\circledast}F_{2}$ 一致。类似地，考虑映射 $m_{Y}:(K\backslash G)\times Y\to\mathcal{B}$；置 $F\widetilde{\circledast}T:=m_{Y\cdot}(F_{K\backslash G}\boxtimes T)$。$(K\backslash G)\times Y$ 上显然的 $K$-作用给出该 $\Omega$-复形上的 $K_{\Omega}$-作用，其不变量与 $F\overset{c}{\circledast}T$ 一致。

% [Tutor Note: 原书 p.264 末句「coinside」拼写如此（应为 coincide），底本如此。]

\bdp{7.6.7} 把 $Y$ 上弱 $K$-等变 $\Omega$-复形的范畴记为 $C(K^{\backslash}_{\backslash}Y,\Omega)$，相应的同伦范畴与 $\mathcal{D}$-导出范畴记为 $K(K^{\backslash}_{\backslash}Y,\Omega)$、$D(K^{\backslash}_{\backslash}Y,\Omega)$（弱等变 $\Omega$-复形的态射若作为普通 $\Omega$-复形是 $\mathcal{D}$-拟同构，则称为 $\mathcal{D}$-拟同构）。

\bdp{7.6.8} \emph{注.} (i) 遗忘函子 $C(\mathcal{B},\Omega)\to C(K^{\backslash}_{\backslash}Y,\Omega)$ 有左右伴随函子 $c^{l},c^{r}:C(K^{\backslash}_{\backslash}Y,\Omega)\to C(\mathcal{B},\Omega)$，$c^{l}(F_{Y})=U(\mathfrak{k}_{\Omega})\otimes_{U(\mathfrak{k})}F_{Y}$，$c^{r}(F_{Y})=\mathrm{Hom}_{U(\mathfrak{k})}(U(\mathfrak{k}_{\Omega}),F_{Y})$。这些函子保持拟同构，故给出导出范畴之间的伴随函子。

(ii) 遗忘函子 $C(K^{\backslash}_{\backslash}Y,\Omega)\to C(Y,\Omega)$ 有右伴随函子 $\mathrm{Ind}:C(Y,\Omega)\to C(K^{\backslash}_{\backslash}Y,\Omega)$，$\mathrm{Ind}(T^{\cdot}_{Y})=p_{*}m^{*}(T^{\cdot}_{Y})$，其中 $m,p:K\times Y\rightrightarrows Y$ 是作用映射与投影映射。这些函子保持拟同构，故给出导出范畴之间的伴随函子。复合 $c^{r}\mathrm{Ind}$ 是投影 $Y\to\mathcal{B}$ 的前推函子。

(iii) 注 7.6.6 (ii) 对弱等变 $\Omega$-复形仍然有效。

% [Tutor Note: 原书 p.264 引「Remark 7.6.6 (ii)」，但 7.6.6 并无 (ii)，底本如此，译文依原样。]

(iv) 设 $f:Z\to Y$ 是带 $K$-作用的光滑簇的态射。7.3.6 中正像函子的构造无需改动即可搬到弱等变设定，故有函子 $f_{*}=Rf_{\cdot}:D(K^{\backslash}_{\backslash}Z,\Omega)\to D(K^{\backslash}_{\backslash}Y,\Omega)$。函子 $f_{*}$ 与\SourcePage{265}(i)、(ii) 中所述的函子交换。对 7.2.8、7.3.6 中的拉回函子 $f^{\cdot}_{\Omega}$ 也一样。

(v) 下面是 7.6.1 的弱等变版本。设 7.6.1 中的 $Y$ 还带有仿射代数群 $G'$ 的作用，且与 $G$-作用交换（我们把它写成右作用）。考虑 $Y$ 上带有交换的 $K_{\Omega}$-作用与 $G$-作用的 $\Omega$-复形范畴 $C(K^{\backslash}_{\backslash}Y_{/}^{/}G',\Omega)=C(\mathcal{B}_{/}^{/}G',\Omega)$。则相应的导出范畴 $D(\mathcal{B}_{/}^{/}G',\Omega)$ 是 $\mathcal{H}$-模。$\mathcal{H}$-作用的定义方式与 7.6.1 相同。注 7.6.6 仍然有效。

% [Tutor Note: 原书 p.265 此处作「commuting K_Ω- and G-actions」，按上下文应为 G′-作用；底本如此，译文依原样作 G-作用。]

\bdp{7.6.9} 下面描述上述等变范畴的 $\mathcal{D}$-模对应物（细节见 [BL]）。对 $Y$ 上的 $\mathcal{D}$-模 $M$，$M$ 上的\emph{弱 $K$-作用} 是把 $M$ 看作 $\mathcal{O}_{Y}$-模时的 $K$-作用，使得对任意 $k\in K$，平移 $k^{*}M\xrightarrow{\sim}M$ 是 $\mathcal{D}$-模的态射。带弱 $K$-作用的 $\mathcal{D}$-模称为\emph{弱 $K$-等变} $\mathcal{D}$-模；其范畴记为 $\mathcal{M}(K^{\backslash}_{\backslash}Y)$（按惯例用 $\mathcal{M}^{\ell}$ 或 $\mathcal{M}^{r}$ 指明左、右 $\mathcal{D}$-模）。记号 $C(K^{\backslash}_{\backslash}Y,\mathcal{D})$、$K(K^{\backslash}_{\backslash}Y,\mathcal{D})$、$D(K^{\backslash}_{\backslash}Y,\mathcal{D})=D(K^{\backslash}_{\backslash}Y)$ 的含义是清楚的（参看 7.2）。

7.2.2 中的函子 $\mathcal{D}$ 与 $\Omega$ 把弱等变复形送到弱等变复形，于是有伴随 DG 函子
\begin{equation}\tag{320}
\mathcal{D}:C(K^{\backslash}_{\backslash}Y,\Omega)\to C(K^{\backslash}_{\backslash}Y,\mathcal{D}),\qquad
\Omega:C(K^{\backslash}_{\backslash}Y,\mathcal{D})\to C(K^{\backslash}_{\backslash}Y,\Omega),
\end{equation}
以及三角范畴之间互逆的等价
\begin{equation}\tag{321}
D(K^{\backslash}_{\backslash}Y,\mathcal{D})\leftrightarrows D(K^{\backslash}_{\backslash}Y,\Omega).
\end{equation}

按惯例，把这些等同起来的范畴记为 $D(K^{\backslash}_{\backslash}Y)$。

\bdp{7.6.10} \emph{注.} 对弱 $K$-等变 $\mathcal{D}$-模 $M$，$Y$ 上的 $\mathfrak{k}$-作用以两种方式提升到 $\mathcal{O}$-模 $M$：或者是由 $M$ 上的 $K$-作用定义的无穷小作用，或者经由 $Y$ 上的 $\mathfrak{k}$-作用 $\sigma:\mathfrak{k}\to\Theta_{Y}$ 与 $M$ 上的 $\mathcal{D}$-模结构。分别记这两个作用为 $\xi,m\mapsto\mathrm{Lie}_{\xi}m$ 与 $\sigma_{\xi}m$。置 $\xi^{\natural}m:=\mathrm{Lie}_{\xi}m-\sigma_{\xi}m$。则 $\xi^{\natural}\in\mathrm{End}_{\mathcal{D}}M$，并且

% 尾部接 p.266（原书 p.265 末句未完）。
